%2multibyte Version: 5.50.0.2953 CodePage: 1252

\documentclass[notes=show]{beamer}
%%%%%%%%%%%%%%%%%%%%%%%%%%%%%%%%%%%%%%%%%%%%%%%%%%%%%%%%%%%%%%%%%%%%%%%%%%%%%%%%%%%%%%%%%%%%%%%%%%%%%%%%%%%%%%%%%%%%%%%%%%%%%%%%%%%%%%%%%%%%%%%%%%%%%%%%%%%%%%%%%%%%%%%%%%%%%%%%%%%%%%%%%%%%%%%%%%%%%%%%%%%%%%%%%%%%%%%%%%%%%%%%%%%%%%%%%%%%%%%%%%%%%%%%%%%%
\usepackage{amsfonts}
\usepackage{graphicx}
\usepackage{amssymb}
\usepackage{amsmath}
\usepackage{mathpazo}
\usepackage{hyperref}
\usepackage{multimedia}

\setcounter{MaxMatrixCols}{10}
%TCIDATA{OutputFilter=LATEX.DLL}
%TCIDATA{Version=5.50.0.2953}
%TCIDATA{Codepage=1252}
%TCIDATA{<META NAME="SaveForMode" CONTENT="2">}
%TCIDATA{BibliographyScheme=Manual}
%TCIDATA{Created=Wednesday, May 01, 2013 14:54:39}
%TCIDATA{LastRevised=Monday, November 15, 2021 13:59:12}
%TCIDATA{<META NAME="GraphicsSave" CONTENT="32">}
%TCIDATA{<META NAME="DocumentShell" CONTENT="Other Documents\SW\Slides - Beamer">}
%TCIDATA{Language=American English}
%TCIDATA{CSTFile=beamer.cst}
%TCIDATA{PageSetup=72,72,72,83,0}
%TCIDATA{Counters=arabic,1}
%TCIDATA{AllPages=
%H=36
%F=36
%}


\newenvironment{stepenumerate}{\begin{enumerate}[<+->]}{\end{enumerate}}
\newenvironment{stepitemize}{\begin{itemize}[<+->]}{\end{itemize} }
\newenvironment{stepenumeratewithalert}{\begin{enumerate}[<+-| alert@+>]}{\end{enumerate}}
\newenvironment{stepitemizewithalert}{\begin{itemize}[<+-| alert@+>]}{\end{itemize} }
\input{tcilatex}
\usetheme{JuanLesPins}
\usecolortheme{whale}

\begin{document}

\title{Lecture 5: \\
Generalised Least Squares}
\author{Katerina Petrova \ \ \ \ \ \ \ \ \ \ \ \ \ \ \ \ \ \ \ \ \ \ \ \ \ \
\ \ \ \ \ \ \ \ \ \ \ \ \ \ \ \ \ \ \ \ \ \ \ \ \ \ \ \ \ \ \ \ \ \ \ \ \ \
\ \ \ \ \ \ \ \ \ \ \ \ \ \ \ \ \ \ \ \ \ \ \ }
\date{}
\maketitle

\section{The Model}

%TCIMACRO{\TeXButton{BeginFrame}{\begin{frame}}}%
%BeginExpansion
\begin{frame}%
%EndExpansion

\QTR{frametitle}{Classical Assumptions for OLS}

\begin{itemize}
\item We have been working with the following model 
\begin{equation}
y=X\beta +\varepsilon  \label{model}
\end{equation}

\item Let's recall the assumptions we made
\end{itemize}

\begin{enumerate}
\item The model in (\ref{model}) is correctly specified, and crucially,
linear in $\beta $ and $\varepsilon .$

\item The columns of $X$ are linearly independent, so that $rank(X^{\prime
}X)=rank(X)=k<n.$

\item $\varepsilon _{i}$ are i.i.d. variables and $X$ are exogenous
identically distributed regressors, such that $\mathbb{E}[\varepsilon
_{i}|X]=0$ and $\mathbb{E}[\varepsilon _{i}^{2}|X]=\sigma ^{2}.$

\item $p\lim_{n\rightarrow \infty }(\frac{1}{n}X^{\prime
}X)=p\lim_{n\rightarrow \infty }\frac{1}{n}\sum_{i=1}^{n}\mathbf{x}_{i}%
\mathbf{x}_{i}^{\prime }=Q,$ where $\left\Vert Q\right\Vert <\infty $ and $Q$
is positive definite matrix.
\end{enumerate}

%TCIMACRO{\TeXButton{EndFrame}{\end{frame}}}%
%BeginExpansion
\end{frame}%
%EndExpansion

\section{Consequences for OLS}

%TCIMACRO{\TeXButton{BeginFrame}{\begin{frame}}}%
%BeginExpansion
\begin{frame}%
%EndExpansion

\QTR{frametitle}{Relaxing Assumption 3}

\begin{itemize}
\item In this lecture, we will maintain all the classical assumptions above
but we will relax the assumption of no heteroskedasticity, no serial
correlation in the error term

\item In particular, we continue to assume that $\varepsilon $ is a random
vector and $X$ are exogenous regressors so that $\mathbb{E}[\varepsilon
|X]=0 $

\item But now, we will relax the assumption of i.i.d. $\varepsilon $, so
that $\mathbb{E}[\varepsilon \varepsilon ^{\prime }|X]=\left[ 
\begin{array}{cccc}
var(\varepsilon _{1}) & cov(\varepsilon _{1},\varepsilon _{2}) & \cdots & 
cov(\varepsilon _{1},\varepsilon _{n}) \\ 
cov(\varepsilon _{2},\varepsilon _{1}) & var(\varepsilon _{2}) &  &  \\ 
\vdots &  & \ddots &  \\ 
cov(\varepsilon _{n},\varepsilon _{1}) &  &  & var(\varepsilon _{n})%
\end{array}%
\right] =\sigma ^{2}\Omega ,$ where $\Omega $ can be a full p.d. covariance
matrix.

\item Assumption 3: $\varepsilon _{i}$ are random variables and $X$ are
exogenous regressors, satisfying $\mathbb{E}[\varepsilon _{i}|X]=0$ and $%
\mathbb{E}[\varepsilon \varepsilon ^{\prime }|X]=\sigma ^{2}\Omega $
\end{itemize}

%TCIMACRO{\TeXButton{EndFrame}{\end{frame}}}%
%BeginExpansion
\end{frame}%
%EndExpansion

%TCIMACRO{\TeXButton{BeginFrame}{\begin{frame}}}%
%BeginExpansion
\begin{frame}%
%EndExpansion

\QTR{frametitle}{What has changed?}

$\widehat{\beta }_{n}^{OLS}$ is still unbiased. Recall that $\widehat{\beta }%
_{n}^{OLS}=\beta +(X^{\prime }X)^{-1}X^{\prime }\varepsilon .$

\begin{eqnarray*}
\mathbb{E}[\widehat{\beta }_{n}^{OLS}] &=&\beta +\mathbb{E}[(X^{\prime
}X)^{-1}X^{\prime }\varepsilon ] \\
&&=\mathbb{E[E[}(X^{\prime }X)^{-1}X^{\prime }\varepsilon |X]] \\
&=&\beta +\mathbb{E[}(X^{\prime }X)^{-1}X^{\prime }\underset{=0}{\underbrace{%
\mathbb{E[}\varepsilon |X]}]}=\beta
\end{eqnarray*}

%TCIMACRO{\TeXButton{EndFrame}{\end{frame}}}%
%BeginExpansion
\end{frame}%
%EndExpansion

%TCIMACRO{\TeXButton{BeginFrame}{\begin{frame}}}%
%BeginExpansion
\begin{frame}%
%EndExpansion

\QTR{frametitle}{What has changed?}

\begin{itemize}
\item Crucially, it is also still consistent.
\end{itemize}

\begin{eqnarray*}
p\lim_{n\rightarrow \infty }\widehat{\beta }_{n}^{OLS} &=&\beta
+p\lim_{n\rightarrow \infty }\left[ \left( \frac{1}{n}\tsum%
\nolimits_{i=1}^{n}\mathbf{x}_{i}\mathbf{x}_{i}^{\prime }\right) ^{-1}\frac{1%
}{n}\tsum\nolimits_{i=1}^{n}\mathbf{x}_{i}\varepsilon _{i}\right] \\
&&\overset{\text{CMT}}{=}\beta +\left( p\lim_{n\rightarrow \infty }\frac{1}{n%
}\tsum\nolimits_{i=1}^{n}\mathbf{x}_{i}\mathbf{x}_{i}^{\prime }\right)
^{-1}p\lim_{n\rightarrow \infty }\left( \frac{1}{n}\tsum\nolimits_{i=1}^{n}%
\mathbf{x}_{i}\varepsilon _{i}\right) \\
&&\overset{\text{WLLN}}{=}\beta +Q^{-1}\underset{=0}{\underbrace{\mathbb{E[}%
\mathbf{x}_{i}\varepsilon _{i}]}} \\
&=&\beta
\end{eqnarray*}

%TCIMACRO{\TeXButton{EndFrame}{\end{frame}}}%
%BeginExpansion
\end{frame}%
%EndExpansion

%TCIMACRO{\TeXButton{BeginFrame}{\begin{frame}}}%
%BeginExpansion
\begin{frame}%
%EndExpansion

\QTR{frametitle}{What has changed?}

\begin{gather*}
var[\widehat{\beta }_{n}^{OLS}]=\mathbb{E}[(\widehat{\beta }_{n}^{OLS}-\beta
)(\widehat{\beta }_{n}^{OLS}-\beta )^{\prime }] \\
=\mathbb{E}[((X^{\prime }X)^{-1}X^{\prime }\varepsilon )((X^{\prime
}X)^{-1}X^{\prime }\varepsilon )^{\prime }] \\
\overset{\text{LIE, CT}}{=}\mathbb{E}[((X^{\prime }X)^{-1}X^{\prime }\mathbb{%
E[}\varepsilon \varepsilon ^{\prime }|X]X(X^{\prime }X)^{-1}] \\
=\sigma ^{2}\mathbb{E}[((X^{\prime }X)^{-1}X^{\prime }\Omega X(X^{\prime
}X)^{-1}].
\end{gather*}

%TCIMACRO{\TeXButton{EndFrame}{\end{frame}}}%
%BeginExpansion
\end{frame}%
%EndExpansion

%TCIMACRO{\TeXButton{BeginFrame}{\begin{frame}}}%
%BeginExpansion
\begin{frame}%
%EndExpansion

\QTR{frametitle}{What has changed?}

\begin{gather*}
\sqrt{n}(\widehat{\beta }_{n}^{OLS}-\beta )=\left( \frac{1}{n}%
\tsum\nolimits_{i=1}^{n}\mathbf{x}_{i}\mathbf{x}_{i}^{\prime }\right) ^{-1}%
\frac{1}{\sqrt{n}}\tsum\nolimits_{i=1}^{n}\mathbf{x}_{i}\varepsilon _{i} \\
=\underset{\rightarrow _{p}Q^{-1}}{\underbrace{\left( \frac{1}{n}%
\tsum\nolimits_{i=1}^{n}\mathbf{x}_{i}\mathbf{x}_{i}^{\prime }\right) ^{-1}}}%
\underset{\rightarrow _{d}\mathcal{N}(0,S)}{\underbrace{\frac{1}{\sqrt{n}}%
\tsum\nolimits_{i=1}^{n}\mathbf{x}_{i}\varepsilon _{i}}} \\
\rightarrow _{d}\mathcal{N}(0,Q^{-1}SQ^{-1})\text{ where }S=\limfunc{plim}%
_{n\rightarrow \infty }\frac{1}{n}\tsum\nolimits_{i=1}^{n}\varepsilon
_{i}^{2}\mathbf{x}_{i}\mathbf{x}_{i}^{\prime }
\end{gather*}%
where a non-i.i.d. CLT\ has to be used (see textbook for technical
conditions).

%TCIMACRO{\TeXButton{EndFrame}{\end{frame}}}%
%BeginExpansion
\end{frame}%
%EndExpansion

%TCIMACRO{\TeXButton{BeginFrame}{\begin{frame}}}%
%BeginExpansion
\begin{frame}%
%EndExpansion

\QTR{frametitle}{What has changed?}

\begin{itemize}
\item So ignoring the heteroscedasticity/serial correlation of the errors,
we remain unbiased and consistent

\item However, in this case, if we applied the standard OLS we will end up
using the wrong variance

\item In this case, $\widehat{\beta }_{n}^{OLS}$ will no longer be
efficient: we can prove that there exists another linear unbiased estimator
which has smaller variance.

\item Moreover, the variance in the asymptotic distribution is different
(sandwich), so ignoring the heteroskedasticity/serial correlation would
invalidate testing!!!

\item Recall that because $\Omega $ is symmetric positive definite matrix,
so is $\Omega ^{-1}$, and there exists a nonsingular $P$ such that $\Omega
^{-1}=P^{\prime }P$ and equivalently $\Omega =(P^{^{\prime
}}P)^{-1}=P^{-1}P^{\prime -1}.$
\end{itemize}

%TCIMACRO{\TeXButton{EndFrame}{\end{frame}}}%
%BeginExpansion
\end{frame}%
%EndExpansion

\section{The GLS}

%TCIMACRO{\TeXButton{BeginFrame}{\begin{frame}}}%
%BeginExpansion
\begin{frame}%
%EndExpansion

\QTR{frametitle}{Transforming the Model}

\begin{itemize}
\item Let's premultiply our model with $P:$ 
\begin{equation}
\underset{\widetilde{y}}{\underbrace{Py}}=\underset{\widetilde{X}}{%
\underbrace{PX}}\beta +\underset{\widetilde{\varepsilon }}{\underbrace{%
P\varepsilon }}  \label{transformed}
\end{equation}

\item Then, $var(\widetilde{\varepsilon })=Pvar(\varepsilon )P^{\prime
}=P\sigma ^{2}\Omega P^{\prime }=\sigma ^{2}PP^{-1}P^{\prime -1}P^{\prime
}=\sigma ^{2}I_{n}.$

\item So by premultiplying the model with $P$, we obtained our original
assumption 3.

\item This is known as standardising (since premultiplying by $P$ is the
matrix equivalent to deviding by the standard deviation)
\end{itemize}

%TCIMACRO{\TeXButton{EndFrame}{\end{frame}}}%
%BeginExpansion
\end{frame}%
%EndExpansion

%TCIMACRO{\TeXButton{BeginFrame}{\begin{frame}}}%
%BeginExpansion
\begin{frame}%
%EndExpansion

\QTR{frametitle}{The GLS estimator}

\begin{itemize}
\item The transformation above gives rise to a new estimator

\item $\arg \min \left\Vert \widetilde{y}-\widetilde{X}\beta \right\Vert
^{2}=(\widetilde{y}-\widetilde{X}\beta )^{\prime }(\widetilde{y}-\widetilde{X%
}\beta )$

$\Rightarrow \widehat{\beta }^{GLS}=(\widetilde{X}^{\prime }\widetilde{X}%
)^{-1}\widetilde{X}^{\prime }\widetilde{y}=(X^{\prime }P^{\prime
}PX)^{-1}X^{\prime }P^{\prime }Py=(X^{\prime }\Omega ^{-1}X)^{-1}X^{\prime
}\Omega ^{-1}y.$
\end{itemize}

%TCIMACRO{\TeXButton{EndFrame}{\end{frame}}}%
%BeginExpansion
\end{frame}%
%EndExpansion

%TCIMACRO{\TeXButton{BeginFrame}{\begin{frame}}}%
%BeginExpansion
\begin{frame}%
%EndExpansion

\QTR{frametitle}{The GLS estimator}

\begin{itemize}
\item $\widehat{\beta }^{GLS}$ is unbiased: $\widehat{\beta }%
^{GLS}=(X^{\prime }\Omega ^{-1}X)^{-1}X^{\prime }\Omega ^{-1}(X\beta
+\varepsilon )$

\item $\mathbb{E}[\widehat{\beta }^{GLS}]=\beta +\mathbb{E}\left[ (X^{\prime
}\Omega ^{-1}X)^{-1}X^{\prime }\Omega ^{-1}\mathbb{E[}\varepsilon |X]\right]
=\beta .$

\item $var(\widehat{\beta }^{GLS})=\mathbb{E}[(\widehat{\beta }^{GLS}-\beta
)(\widehat{\beta }^{GLS}-\beta )^{\prime }]=\mathbb{E}\left[ ((X^{\prime
}\Omega ^{-1}X)^{-1}X^{\prime }\Omega ^{-1}\varepsilon )((X^{\prime }\Omega
^{-1}X)^{-1}X^{\prime }\Omega ^{-1}\varepsilon )^{\prime }\right] =$

$=\mathbb{E}\left[ ((X^{\prime }\Omega ^{-1}X)^{-1}X^{\prime }\Omega ^{-1}%
\mathbb{E[}\varepsilon \varepsilon ^{\prime }|X]\Omega ^{-1}X(X^{\prime
}\Omega ^{-1}X)^{-1}\right] =\mathbb{\sigma }^{2}\mathbb{E}(X^{\prime
}\Omega ^{-1}X)^{-1}.$
\end{itemize}

%TCIMACRO{\TeXButton{EndFrame}{\end{frame}}}%
%BeginExpansion
\end{frame}%
%EndExpansion

%TCIMACRO{\TeXButton{BeginFrame}{\begin{frame}}}%
%BeginExpansion
\begin{frame}%
%EndExpansion

\QTR{frametitle}{Weighted Least Squares}

\begin{itemize}
\item Consider the special case when $\Omega $ is a diagonal matrix%
\begin{equation*}
\Omega =\left[ 
\begin{array}{ccc}
\omega _{ii} & \cdots & 0 \\ 
\vdots & \ddots & \vdots \\ 
0 & \cdots & \omega _{nn}%
\end{array}%
\right]
\end{equation*}

\item $P$ is a diagonal matrix, containing as elements $\frac{1}{\sqrt{%
\omega _{jj}}},$ $j=1,...,n$

\item The GLS estimator is then given by $\widehat{\beta }^{GLS}=(X^{\prime
}\Omega ^{-1}X)^{-1}X^{\prime }\Omega ^{-1}y=\left(
\tsum\nolimits_{i=1}^{n}\omega _{ii}^{-1}\mathbf{x}_{i}\mathbf{x}%
_{i}^{\prime }\right) ^{-1}\tsum\nolimits_{i=1}^{n}\omega _{ii}^{-1}\mathbf{x%
}_{i}y_{i}$

\item This estimator is often referred to as WLS because of the weighted sums
\end{itemize}

%TCIMACRO{\TeXButton{EndFrame}{\end{frame}}}%
%BeginExpansion
\end{frame}%
%EndExpansion

\section{Gauss Markov Theorem}

%TCIMACRO{\TeXButton{BeginFrame}{\begin{frame}}}%
%BeginExpansion
\begin{frame}%
%EndExpansion

\QTR{frametitle}{Efficiency}

\begin{itemize}
\item With heteroskedasticity/autocorrelation, OLS - not efficient

\item In fact GLS has smaller variance in this case

\item To see this, recall that $\Omega =P^{-1}P^{\prime -1}$, compute
\end{itemize}

\newline
\begin{eqnarray*}
&&var(\widehat{\beta }^{OLS}|X)-var(\widehat{\beta }^{GLS}|X) \\
&=&\sigma ^{2}\left( X^{\prime }X\right) ^{-1}X^{\prime }\Omega X\left(
X^{\prime }X\right) ^{-1}-\sigma ^{2}\left( X^{\prime }\Omega ^{-1}X\right)
^{-1} \\
&=&\sigma ^{2}\left( X^{\prime }X\right) ^{-1}X^{\prime
}P^{-1}(I_{n}-PX\left( X^{\prime }\Omega ^{-1}X\right) ^{-1}X^{\prime
}P^{\prime })P^{\prime -1}X\left( X^{\prime }X\right) ^{-1} \\
&\geq &0
\end{eqnarray*}%
since $\left[ I_{n}-PX\left( X^{\prime }\Omega ^{-1}X\right) ^{-1}X^{\prime
}P^{\prime }\right] $ is symmetric idempotent, and hence positive
semi-definite.

%TCIMACRO{\TeXButton{EndFrame}{\end{frame}}}%
%BeginExpansion
\end{frame}%
%EndExpansion

%TCIMACRO{\TeXButton{BeginFrame}{\begin{frame}}}%
%BeginExpansion
\begin{frame}%
%EndExpansion

\QTR{frametitle}{Gauss Markov Theorem for GLS}

\begin{theorem}[Gauss Markov]
Let $\widetilde{\beta }$ be any linear unbiased estimator for $\beta $ in
the transformed model (\ref{transformed}). Then all the classical
assumptions hold and $var(\widetilde{\beta })\geq var(\widehat{\beta }%
^{GLS}).$
\end{theorem}

%TCIMACRO{\TeXButton{EndFrame}{\end{frame}}}%
%BeginExpansion
\end{frame}%
%EndExpansion

%TCIMACRO{\TeXButton{BeginFrame}{\begin{frame}}}%
%BeginExpansion
\begin{frame}%
%EndExpansion

\QTR{frametitle}{Gauss Markov Theorem for GLS}

\begin{proof}
The model 
\begin{equation}
\widetilde{y}=\widetilde{X}\beta +\widetilde{\varepsilon }  \label{star}
\end{equation}%
satisfies the classical assumption for OLS, hence the OLS estimator in (\ref%
{star}) which is $\widehat{\beta }^{GLS}$ is BLUE by construction.
\end{proof}

%TCIMACRO{\TeXButton{EndFrame}{\end{frame}}}%
%BeginExpansion
\end{frame}%
%EndExpansion

%TCIMACRO{\TeXButton{BeginFrame}{\begin{frame}}}%
%BeginExpansion
\begin{frame}%
%EndExpansion

\begin{proof}[Alternative Proof]
Consider another linear estimator $\widetilde{\beta }=A^{\prime }y$ where $A$
is an $n\times k$ function of $X.$ Since $\widetilde{\beta }$ is unbiased, 
\begin{equation*}
\beta =\mathbb{E[}\widetilde{\beta }|X]=A^{\prime }\mathbb{E}\left[ y|X%
\right] =A^{\prime }\mathbb{E}\left[ \left( X\beta +\varepsilon \right) |X%
\right] =A^{\prime }X\beta ,
\end{equation*}%
hence $A^{\prime }X=I_{k}.$ Moreover, 
\begin{eqnarray*}
var\left( \widetilde{\beta }|X\right)  &=&var\left( A^{\prime }y|X\right)
=A^{\prime }var\left( y|X\right) A \\
&=&A^{\prime }var\left( \varepsilon |X\right) A=\sigma ^{2}A^{\prime }\Omega
A.
\end{eqnarray*}%
We also know, that $var\left( \hat{\beta}^{GLS}|X\right) =\sigma
^{2}(X^{\prime }\Omega ^{-1}X)^{-1}.$ Now, let's show that $var\left( 
\widetilde{\beta }|X\right) \geq var\left( \hat{\beta}^{GLS}|X\right) ,$
that is $A^{\prime }\Omega A-(X^{\prime }\Omega ^{-1}X)^{-1}$ is positive
semi-definite matrix.
\end{proof}

%TCIMACRO{\TeXButton{EndFrame}{\end{frame}}}%
%BeginExpansion
\end{frame}%
%EndExpansion

%TCIMACRO{\TeXButton{BeginFrame}{\begin{frame}}}%
%BeginExpansion
\begin{frame}%
%EndExpansion

\begin{proof}[Alternative Proof]
Define $C=P^{\prime -1}A-PX\left( X^{\prime }\Omega ^{-1}X\right) ^{-1},$
noting that 
\begin{gather*}
C^{\prime }PX\left( X^{\prime }\Omega ^{-1}X\right) ^{-1} \\
=\left( P^{^{\prime }-1}A-PX\left( X^{\prime }\Omega ^{-1}X\right)
^{-1}\right) ^{\prime }PX\left( X^{\prime }\Omega ^{-1}X\right) ^{-1} \\
=A^{\prime }P^{-1}PX\left( X^{\prime }\Omega ^{-1}X\right) ^{-1}-\left(
X^{\prime }\Omega ^{-1}X\right) ^{-1}X^{\prime }P^{\prime }PX\left(
X^{\prime }\Omega ^{-1}X\right) ^{-1} \\
=0
\end{gather*}%
since $A^{\prime }X=I_{k}.$
\end{proof}

%TCIMACRO{\TeXButton{EndFrame}{\end{frame}}}%
%BeginExpansion
\end{frame}%
%EndExpansion

%TCIMACRO{\TeXButton{BeginFrame}{\begin{frame}}}%
%BeginExpansion
\begin{frame}%
%EndExpansion

\begin{proof}[Alternative Proof]
We have that 
\begin{gather*}
A^{\prime }\Omega A-(X^{\prime }\Omega ^{-1}X)^{-1}=A^{\prime
}P^{-1}P^{\prime -1}A-(X^{\prime }\Omega ^{-1}X)^{-1} \\
=\left( C+PX\left( X^{\prime }\Omega ^{-1}X\right) ^{-1}\right) ^{\prime
}\left( C+PX\left( X^{\prime }\Omega ^{-1}X\right) ^{-1}\right) -(X^{\prime
}\Omega ^{-1}X)^{-1} \\
=C^{\prime }C+C^{\prime }PX\left( X^{\prime }\Omega ^{-1}X\right)
^{-1}+\left( X^{\prime }\Omega ^{-1}X\right) ^{-1}X^{\prime }P^{\prime }C+ \\
\left( X^{\prime }\Omega ^{-1}X\right) ^{-1}X^{\prime }\Omega ^{-1}X\left(
X^{\prime }\Omega ^{-1}X\right) ^{-1}-(X^{\prime }\Omega ^{-1}X)^{-1} \\
=C^{\prime }C\geq 0
\end{gather*}%
since $C^{\prime }C$ is a quadratic form, hence positive semi-definite. So, $%
\widehat{\beta }^{OLS}$ is efficient among a class of linear estimators.
\end{proof}

%TCIMACRO{\TeXButton{EndFrame}{\end{frame}}}%
%BeginExpansion
\end{frame}%
%EndExpansion

%TCIMACRO{\TeXButton{BeginFrame}{\begin{frame}}}%
%BeginExpansion
\begin{frame}%
%EndExpansion

\begin{proof}[Alternative Proof 2]
\begin{gather*}
A^{\prime }\Omega A-(X^{\prime }\Omega ^{-1}X)^{-1}=A^{\prime
}P^{-1}P^{\prime -1}A-A^{\prime }X(X^{\prime }\Omega ^{-1}X)^{-1}X^{\prime }A
\\
=A^{\prime }P^{-1}\left( I_{n}-PX(X^{\prime }\Omega ^{-1}X)^{-1}X^{\prime
}P^{\prime }\right) P^{\prime -1}A\geq 0
\end{gather*}%
since $(I_{n}-PX\left( X^{\prime }\Omega ^{-1}X\right) ^{-1}X^{\prime
}P^{\prime })$ is idempotent and hence positive semi-definite.
\end{proof}

%TCIMACRO{\TeXButton{EndFrame}{\end{frame}}}%
%BeginExpansion
\end{frame}%
%EndExpansion

\section{More properties}

%TCIMACRO{\TeXButton{BeginFrame}{\begin{frame}}}%
%BeginExpansion
\begin{frame}%
%EndExpansion

\QTR{frametitle}{Equivalence of OLS and GLS}

\begin{itemize}
\item Note that whenever $\Omega =I_{n},$ so that $\mathbb{E}[\varepsilon
\varepsilon ^{\prime }|X]=\sigma ^{2}\Omega =\sigma ^{2}I_{n},$ $\widehat{%
\beta }^{GLS}=(X^{\prime }\Omega ^{-1}X)^{-1}X^{\prime }\Omega
^{-1}y=(X^{\prime }X)^{-1}X^{\prime }y=\widehat{\beta }^{OLS}$

\item So the GLS estimator reduces to OLS
\end{itemize}

%TCIMACRO{\TeXButton{EndFrame}{\end{frame}}}%
%BeginExpansion
\end{frame}%
%EndExpansion

%TCIMACRO{\TeXButton{BeginFrame}{\begin{frame}}}%
%BeginExpansion
\begin{frame}%
%EndExpansion

\QTR{frametitle}{Projection Interpretation}

\begin{itemize}
\item The matrix $R_{X}=X(X^{\prime }\Omega ^{-1}X)^{-1}X^{\prime }\Omega
^{-1}$ is idempotent, but no longer symmetric

\item This implies that $R_{X}$ is oblique (but not orthogonal) projection
of $y$ on the space spanned by $X$

\item It also implies that the GLS resudual $\widehat{\varepsilon }^{GLS}$
is not orthogonal to the space spanned by $X$

\item In fact, $\widehat{\varepsilon }^{GLS}$ is orthogonal to $\Omega
^{-1}X.$

\item Another way to think about the GLS procedure is as the minimiser of
the Mahalanobis norm%
\begin{gather*}
\arg \min_{\beta }\left\Vert y-X\beta \right\Vert _{V}^{2}=\arg \min_{\beta
}(y-X\beta )^{\prime }[\Omega ^{-1}/\sigma ^{2}](y-X\beta ) \\
=\arg \min_{\beta }(y-X\beta )^{\prime }P^{\prime }P(y-X\beta )\text{ 
{\small since} }{\small \sigma }^{2}\text{ {\small doesn't matter}} \\
=\arg \min_{\beta }(\widetilde{y}-\widetilde{X}\beta )^{\prime }(\widetilde{y%
}-\widetilde{X}\beta )
\end{gather*}
\end{itemize}

%TCIMACRO{\TeXButton{EndFrame}{\end{frame}}}%
%BeginExpansion
\end{frame}%
%EndExpansion

%TCIMACRO{\TeXButton{BeginFrame}{\begin{frame}}}%
%BeginExpansion
\begin{frame}%
%EndExpansion

\QTR{frametitle}{Inference}

\begin{itemize}
\item We can use the asymptotic distribution of GLS to test linear
restrictions
\end{itemize}

\begin{theorem}
Under the null hypothesis $R\beta ^{0}=c$,%
\begin{equation*}
\frac{(R\widehat{\beta }^{GLS}-c)^{\prime }[R(X^{\prime }\Omega
^{-1}X)^{-1}R^{\prime }]^{-1}(R\widehat{\beta }^{GLS}-c)}{\hat{\sigma}^{2}}%
\sim \chi ^{2}(r).
\end{equation*}
\end{theorem}

\begin{itemize}
\item The left-hand side above can serve as a test statistic for the linear
hypothesis $R\beta ^{0}=c.$

\item The $z$ test is a special case of this when $r=1$

\item Note for this result to hold, we treat $\Omega $ as known
\end{itemize}

%TCIMACRO{\TeXButton{EndFrame}{\end{frame}}}%
%BeginExpansion
\end{frame}%
%EndExpansion

\section{Feasible and Infeasible GLS}

%TCIMACRO{\TeXButton{BeginFrame}{\begin{frame}}}%
%BeginExpansion
\begin{frame}%
%EndExpansion

\QTR{frametitle}{A serious drawback: Feasible GLS}

\begin{itemize}
\item In practice, $\Omega $ is typically unknown so that the GLS estimator
is not available.

\item Substituting an estimator $\widehat{\Omega }$ for $\Omega $ yields the 
\textbf{feasible} generalized least squares (FGLS) estimator%
\begin{equation*}
\widehat{\beta }^{FGLS}=(X^{\prime }\widehat{\Omega }^{-1}X)^{-1}X^{\prime }%
\widehat{\Omega }^{-1}y
\end{equation*}

\item Note, however, that $\Omega $ contains too many $n(n+1)/2$ parameters.

\item Proper estimation of $\Omega $ would not be possible unless further
structure on the elements of $\Omega $ are imposed.

\item Under different assumptions on $var(\varepsilon ),$ $\Omega $ has a
simpler structure with much fewer (say, $p\ll n$) unknown parameters and may
be properly estimated.
\end{itemize}

%TCIMACRO{\TeXButton{EndFrame}{\end{frame}}}%
%BeginExpansion
\end{frame}%
%EndExpansion

%TCIMACRO{\TeXButton{BeginFrame}{\begin{frame}}}%
%BeginExpansion
\begin{frame}%
%EndExpansion

\QTR{frametitle}{Feasible GLS}

\begin{itemize}
\item The properties of FGLS estimation crucially depend on these assumptions

\item A clear disadvantage of the FGLS estimator is that its finite sample
properties are usually unknown

\item Note that $\widehat{\Omega }_{n}$ is, in general, a function of $y$,
so that $\widehat{\beta }^{FGLS}$ is a complex function of the elements of $%
y $

\item It is therefore difficult, and in most cases not possible, to derive
the finite-sample properties, such as expectation and variance, of $\widehat{%
\beta }^{FGLS}$

\item Consequently, the efficiency gain of an FGLS estimator is not at all
clear

\item Deriving exact tests based on the FGLS estimator is also a difficult
job

\item We must rely on the asymptotic properties of $\widehat{\beta }^{FGLS}$
to draw statistical inferences
\end{itemize}

%TCIMACRO{\TeXButton{EndFrame}{\end{frame}}}%
%BeginExpansion
\end{frame}%
%EndExpansion

\section{Robust OLS}

%TCIMACRO{\TeXButton{BeginFrame}{\begin{frame}}}%
%BeginExpansion
\begin{frame}%
%EndExpansion

\QTR{frametitle}{Heteroskedasticity Robust Standard Errors}

\begin{itemize}
\item Consider the simple case of diagonal $\Omega $

\item Since GLS or WLS are difficult to obtain unless additional assumptions
are imposed in order to estimate $\Omega ,$ is it possible to use OLS
instead?

\item Recall that 
\begin{equation*}
\sqrt{n}(\widehat{\beta }_{n}^{OLS}-\beta )\rightarrow _{d}\mathcal{N}%
(0,Q^{-1}SQ^{-1}),\text{ }S=\limfunc{plim}_{n\rightarrow \infty }\frac{1}{n}%
\tsum\nolimits_{i=1}^{n}\varepsilon _{i}^{2}\mathbf{x}_{i}\mathbf{x}%
_{i}^{\prime }
\end{equation*}

\item In order to do valid inference we must use the correct variance for $%
\widehat{\beta }_{n}^{OLS}$

\item We need a consistent estimator for $Q^{-1}SQ^{-1}$

\item We already know how to estimate $Q,$ using Matching moments $\hat{Q}=%
\frac{1}{n}\tsum\nolimits_{i=1}^{n}\mathbf{x}_{i}\mathbf{x}_{i}^{\prime }.$
\end{itemize}

%TCIMACRO{\TeXButton{EndFrame}{\end{frame}}}%
%BeginExpansion
\end{frame}%
%EndExpansion

%TCIMACRO{\TeXButton{BeginFrame}{\begin{frame}}}%
%BeginExpansion
\begin{frame}%
%EndExpansion

\QTR{frametitle}{Heteroskedasticity Robust Standard Errors}

\begin{itemize}
\item A consistent estimate for $S$ is given by $\frac{1}{n}%
\tsum\nolimits_{i=1}^{n}\hat{\varepsilon}_{i}^{2}\mathbf{x}_{i}\mathbf{x}%
_{i}^{\prime }$ (for which a non-i.i.d. WLLN is required)

\item The resulting robust standard error for $\hat{\beta}_{j}$ is given by $%
SE\left( \hat{\beta}_{j}\right) =$ 
\begin{equation*}
\sqrt{\left[ \frac{\left( \frac{1}{n}\tsum\nolimits_{i=1}^{n}\mathbf{x}_{i}%
\mathbf{x}_{i}^{\prime }\right) ^{-1}\left( \frac{1}{n}\tsum%
\nolimits_{i=1}^{n}\mathbf{x}_{i}\mathbf{x}_{i}^{\prime }\hat{\varepsilon}%
_{i}^{2}\right) \left( \frac{1}{n}\tsum\nolimits_{i=1}^{n}\mathbf{x}_{i}%
\mathbf{x}_{i}^{\prime }\right) ^{-1}}{n}\right] _{jj}}
\end{equation*}

\item These standard errors are known as Heteroskedasticity Consistent
(White) SEs

\item There are also Heteroskedasticity Autocorrelation Consistent (Newey
West) SEs, which allow for autocorrelation in $\varepsilon ,$ so that $%
\Omega $ is not necessarily diagonal

\item This is beyond the scope of the course.
\end{itemize}

%TCIMACRO{\TeXButton{EndFrame}{\end{frame}}}%
%BeginExpansion
\end{frame}%
%EndExpansion

\section{Reading}

%TCIMACRO{\TeXButton{BeginFrame}{\begin{frame}}}%
%BeginExpansion
\begin{frame}%
%EndExpansion

\QTR{frametitle}{Reading}

\begin{itemize}
\item Relevant chapters from the textbook: 4.9, 4.14, 4.15, 7.7.
\end{itemize}

%TCIMACRO{\TeXButton{EndFrame}{\end{frame}}}%
%BeginExpansion
\end{frame}%
%EndExpansion

\end{document}
